% from main.tex
%%%%%%%%%%  scan idx 129  |  folio 116  %%%%%%%%%%
En
particulier, pour que \(F\) soit pseudo-cohérent, il faut et il suffit que \(\Ho^n(F)\) soit cohérent pour tout \(n\).

\emph{Preuve.} L'assertion a) et l'implication \(\Rightarrow\) de b) résultent trivialement de (3.4). L'implication \(\Leftarrow\) de b) découle aussitôt de a) et (2.11.b)).

\textbf{Exemples 3.6.} Soient \((S,A)\) un topos annelé, \(\mathcal C,\mathcal C_0\) les \(S\)-catégories de (1.3.d)) (cf. aussi (2.15)). Pour que les objets de \(\mathcal C_0\) soient cohérents, il faut et il suffit que \(A\) soit cohérent (i.e. précohérent). Voici des exemples où il en est ainsi :
\begin{itemize}[label=--]
\item \(A\) est constant de valeur un anneau noethérien à gauche (SGA 4 IX 2.1).
\item \(S\) est le topos des faisceaux d'ensembles (pour la topologie de Zariski) sur un schéma localement noethérien \(S\)\footnote{Conformément aux nouvelles conventions, nous dirons «~schéma~» pour «~préschéma~», et «~schéma séparé~» pour «~schéma~».}, et \(A=\mathcal O_S\).
\item \(S\) est le topos des faisceaux d'ensembles sur un espace analytique complexe \(S\), \(A=\mathcal O_S\).
\end{itemize}
En revanche, le faisceau des fonctions continues complexes sur un espace topologique n'est en général pas cohérent.

Bien entendu, si \((S,A)\) est défini par un espace annelé, la notion de cohérence envisagée ici coïncide avec la notion usuelle de Serre (EGA \(0_I\) 5.3.1).

\textbf{Proposition 3.7.} \textbf{a)} Soit \(f:(S',A')\to(S,A)\) un morphisme de topos annelés. On suppose \(f\) plat à droite, i.e. que le foncteur \(f^*=A'\otimes_A\cdot\) est exact. Il existe, pour \(F\in\Ob\Dpx(_A S)\), \(G\in\Ob\Dpx^+(_A S)\), une flèche canonique fonctorielle :
\[
  f^*\RHom(F,G)\longrightarrow\RHom(f^*(F),f^*(G)).
\]
